\documentclass[11pt]{article}
\usepackage{fontspec}
\setmainfont{Times New Roman}
\setsansfont{Arial}
\setmonofont{Consolas}
\usepackage{amsmath,amssymb,mathtools}
\usepackage{geometry}
\usepackage{microtype}
\geometry{a4paper,margin=2.35cm}
\setlength{\parindent}{1.25em}
\setlength{\parskip}{0pt}
\makeatletter
\renewcommand\footnotesize{\@setfontsize\footnotesize{7.7}{8.7}\selectfont}
\makeatother
\emergencystretch=100em
\allowdisplaybreaks
\sloppy
\hbadness=10000
\vbadness=10000
\hfuzz=2pt
\vfuzz=10pt
\raggedbottom

\begin{document}

% Унит: Папер31 Сецтион04 ентры 1 в001. Дирецт Интерславиц Латин транслатион фром the Герман финал-аудитед соурце слице, соурце линес 237--243 / сегментс P31-С0053--P31-С0056. Фоотноте цоунтер цонтинуес афтер Папер31 Сецтион03 ентры 5.
\setcounter{footnote}{23}

\subsection*{§4. Матричне представјеньа, следы и дискриминанты при колцах конечного ранга}

В тутом и следујучем параграфу предположеньа о основном колцу будут неколико обчнеје неж доселе, абы објяти такоже порјадкы в числовых и функционалных польах. Дело јест о обчној и једнакој формулацији в суштности знаных фактов.

\emph{Предположенье.} Нехај $\mathfrak R$ буде разширјено колцо области целости $\mathfrak Z$; јединичны елемент колца $\mathfrak R$ уже лежи в $\mathfrak Z$. За кажды $\mathfrak Z$-модул из $\mathfrak R$ -- особливо за само $\mathfrak R$ -- ексистује најманье једна конечна модулна база, составльена из елементов линеарно независных над $\mathfrak Z$. Поред $\alpha_1,\ldots,\alpha_s$, елементы $\beta_1,\ldots,\beta_s$ тогда и само тогда творет линеарно независну модулну базу того самого $\mathfrak Z$-модула, когда детерминант трансформације јест јединица в $\mathfrak Z$.

То предположенье јавно јест исполньено за доселе разсмотрене колца конечного ранга, где $\mathfrak Z$ става польем $P$. Оно такоже јест исполньено за порјадкы в числовых и функционалных польах, где $\mathfrak Z$ јест одповедајуче колцом рационалных целых числ или функционалноју областју.\footnote{Срав. \emph{Идеална теорија}, §3. Елементы из $\mathfrak Z$ означајут се латинскыми буквами.}

\end{document}

